\chapter{Nature's Eternal Return}

\epigraph{``The soul is immortal and has been born many times, and has seen all
things both here and in the underworld.''}{\textcite{plato1997complete}, \emph{Republic} 614d}

\section{The Myth of Er}

The \emph{Republic} ends with a myth. Er, a soldier who is thought to be dead,
revives and tells what he saw in the afterlife. He describes the judgment of
souls and, most importantly, the choice of lives --- the moment when each soul,
before being reborn, chooses the pattern of the life it will live. Some choose
wisely, some foolishly, but all choose according to the character they have
developed in their previous lives.

The myth is not a literal account of the afterlife but a philosophical teaching
about the structure of choice. The soul chooses its life, but the choice is not
arbitrary --- it is conditioned by what the soul has become. The tyrant's soul
chooses a tyrant's life; the philosopher's soul chooses a philosopher's life.
The myth is Plato's way of saying that the character one develops in this life
is the destiny one carries into the next.

\section{The Field of All Possible Actions}

We reframe the Myth of Er as a structural truth about \praxis\ and \physis. The
embedding space of the \physis\ engine --- the high-dimensional vector space in
which every action, every intention, every outcome is represented --- is the
field of all possible actions. Every action the practitioner could take occupies
a point in this space. Every action they have taken occupies another. The
distance between them is the measure of the practitioner's actuality relative
to their potential.

The choice of lives in the myth corresponds to the choice of goals in the
\praxis\ platform. Each practitioner, at each moment, stands before the field
of possible actions and chooses one to actualise. They draw a distinction ---
this, not that --- and register it as a node in the coherence graph. The node
is their trace in the field, their mark on the space of possibilities.

\section{The Judgment of Coherence}

In the myth, souls are judged for their actions and receive their due. In our
framework, the judgment is not administered by a divine judge but by the
structure of nature itself. The coherence score is the judgment: +1 for action
that produced a real effect, 0 for action that left no trace, -1 for action
that violated its own intent.

This judgment is not a punishment or reward. It is a \emph{registration} ---
the recording of what the practitioner's action actually was, as distinct from
what the practitioner intended or believed. The coherence graph is the book of
life, but it is written not by an external scribe but by the practitioner's
own actions, in the very structure of the world.

\section{The Choice of Next Action}

The myth describes souls choosing their next lives. The practitioner, having
reviewed their coherence graph, chooses their next action. The choice is
informed by the record: the practitioner can see which kinds of actions produce
coherence for them and which do not. They can see the patterns of their own past
choices.

But the choice is not determined by the past. The \physis\ engine's dream
function continuously generates novel variations --- proposals for actions that
have never been tried. The practitioner is not condemned to repeat their past
patterns; they can always choose the new, the untried, the unexpected. This is
the freedom that the myth preserves: the soul chooses, and the soul can choose
differently.

\section{Immortality Through Coherence}

The myth concludes with the immortality of the soul. We offer a different kind
of immortality: the persistence of coherent action through the holarchy graph.

Every action that scores +1 establishes a node in the coherence graph that can
serve as a model for future actions --- not only for the practitioner who acted
but for any practitioner who accesses the graph. The action of the practitioner
who trained well, who loved well, who built well --- this action becomes a
node that others can connect to, learn from, build upon.

The \praxis\ platform's community wiki serves this function: the record of what
has worked is shared across practitioners, so that the coherence of one can
inform the action of another. The node does not die with the practitioner; it
persists as long as the graph persists, as long as there are practitioners who
can connect to it. This is immortality not as eternal duration but as eternal
\emph{relevance} --- the action that continues to inform action long after the
actor is gone.

\subsection{Kleos: Immortal Glory Among Mortals}

This is older than it sounds. It is the most archaic of human longings, the one
that moves beneath the \emph{Iliad}: \emph{kleos} --- glory, ``that which is
heard'' --- the undying fame the hero wins when his deed is remembered and
re-sung \parencite[IX]{homer2011iliad}. Achilles chooses a short life with
\emph{kleos} over a long one without it, because the deed that is remembered
outlasts the body that performed it. The Homeric warrior had one technology of
immortality: the song of the bard, which carried the name forward only as long as
there were singers to carry it. The UberWiki is that song made durable. The
coherent deed, deposited as a node, is sung not by a bard who may forget but by a
structure that does not --- and re-sung every time another practitioner connects
to it, builds on it, enacts it again.

There is a deeper logic here, and it is the oldest aristocratic intuition of all:
\emph{agathoi toi agathoi} --- the good to the good, excellence drawn to
excellence. In a commons governed by selection, the most coherent instance of an
action-type accumulates the most connections; it becomes the node others reach for
first; and in becoming so, \emph{the best instance becomes the working definition
of the thing}. The finest act of teaching, of building, of healing that the graph
holds is no longer one teaching among many --- it becomes what teaching, here,
\emph{means}: the reference against which the next is read, the Form that the best
particular has earned the right to be. This is the strange immortality the
practitioner can win. Not that his name is inscribed, but that his deed, having
been the most coherent, \emph{becomes the standard} --- so that everyone who acts
well after him is, in part, enacting him. The \emph{agathos} does not merely
achieve excellence; the excellent deed, deposited, draws the next \emph{agathos}
toward it, and excellence propagates to excellence through the commons. The best
process becomes the definition of the end it aimed at.

So the mortal wins, through the UberWiki, exactly what the Homeric hero wanted and
could secure only by luck of the bards: that his finest action does not pass.
It is taken up into the structure of what others will do, an undying node in a
graph that keeps re-singing it. This is glory among mortals --- not the hollow
fame of the timocratic regime (Book VIII), which sought to be \emph{seen}, but the
true \emph{kleos} of the \emph{agathos}, who sought to act well and, in acting
well, became the form of the act itself. The vir spends his \textsc{ayu} in
coherent deed; the deed becomes the definition; the definition outlives the vir.
That is the whole of it.

\begin{figure}[ht]
\centering
\begin{tikzpicture}[
  mynode/.style={draw, circle, minimum size=6mm, inner sep=1pt, font=\tiny},
  >={Stealth[length=1.5mm]}]
  \foreach \ang in {0,60,120,180,240,300}
    \node[mynode, fill=cohgreen!30, draw=cohgreen] (g\ang) at (\ang:10mm) {};
  \node[mynode, fill=cohgreen!50, draw=cohgreen, minimum size=7mm] (gc) at (0,0) {};

  \foreach \ang in {0,90,180,270}
    \node[mynode, fill=cohgray!30, draw=cohgray] (y\ang) at (\ang:22mm) {};

  \node[mynode, fill=cohred!30, draw=cohred] (r1) at (135:28mm) {};
  \node[mynode, fill=cohred!30, draw=cohred] (r2) at (160:30mm) {};

  \draw[cohgreen!40] (gc) -- (g0) (gc) -- (g60) (gc) -- (g120)
    (gc) -- (g180) (gc) -- (g240) (gc) -- (g300);
  \draw[cohgreen!30] (g0) -- (g60) (g60) -- (g120) (g120) -- (g180)
    (g180) -- (g240) (g240) -- (g300) (g300) -- (g0);
  \draw[cohgray!40] (g0) -- (y0) (g60) -- (y90) (g180) -- (y270) (g240) -- (y180);

  \draw[->, dashed, cohred!60] (r1) to[bend left] ++(10:10mm);
  \draw[->, dashed, cohred!60] (r2) to[bend right] ++(190:8mm);

  \node[font=\tiny, cohgreen, above] at (gc.north) {Coherent (+1)};
  \node[font=\tiny, cohgray, above] at (y0.north) {Inert (0)};
  \node[font=\tiny, cohred, above] at (r1.north) {Failure (--1)};
  \node[font=\tiny, align=center, below=2mm] at (current bounding box.south)
    {Dashed arrows: dream function proposing variations on failed nodes.};
\end{tikzpicture}
\caption{The holarchy graph of actions. Green nodes are coherent actions that
persist; gray nodes are inert; red nodes are failures replayed by the dream
function.\label{fig:holarchy}}
\end{figure}

\section{The Visibility of the They-Self}

Before the closing argument, one preparatory move. The holarchy graph does
something to \emph{das Man} --- Heidegger's anonymous ``they,'' the impersonal
public standard of ``what one does'' that governs inauthentically precisely
because no one authors it and it is everywhere, the air one breathes
\parencite[§27]{heidegger1962being}. The graph \emph{reifies} it: projected into a
shared vector space, the they-self appears as a structure of coherence-weighted
nodes, and what was the invisible horizon of engagement becomes a thing that can
be looked at. Reification, here, is not alienation but its opposite --- for what
can be seen can no longer govern unseen.

The graph makes the they-self visible by giving it the three things it used to
hide behind. A \emph{time}: recurring failure appears as a rhythm, not a run of
accidents, and a rhythm once seen can be broken. A \emph{cause}: every node
carries an edge to the subject who registered it, so the deed ``no one decided but
everyone follows'' is exposed as a chain of abandoned responsibility. And a
\emph{place}: the distance between one's own vector and the statistical centre of
all action is a measurable distance from inauthenticity --- what was everywhere
and nowhere becomes locatable. To see the they-self thus is the first half of
escaping it. The second half --- making it not merely visible but
\emph{answerable} --- is the work of the republic to which we now turn.

\section{The Coherence of the Whole}

The \emph{Republic} ends with an image of cosmic order: the whorl of Necessity
turning the spheres of the cosmos, the Sirens singing each a single note, the
three Fates weaving the destiny of every soul into the fabric of the whole.
Individual choices matter, but they matter within a cosmos that has its own
order, its own music, its own coherence.

The \physis\ engine's holarchy graph is the image of this cosmic order at the
level of action. Every node is connected to others by edges of influence,
causation, and resonance. The practitioner's action is never isolated; it
ripples through the graph, affecting nodes the practitioner may never know
about. The coherence of the whole is not the sum of individual coherences but
the structure of their interconnection --- the shape of the graph itself
\parencite{koestler1967ghost}.

\section{The Bottom-Up Republic}

We have followed Plato's \emph{Republic} through ten books, and at the end we must
confess the one place where we do not follow him at all --- and it is the most
important place. Plato's republic is built from the top down. The philosopher
ascends alone, sees the Good, and \emph{descends to rule}: authority flows
downward from the one who has seen to the many who have not. The city is just
when the wise command and the rest obey. This is the architecture that every
reader since antiquity has found either thrilling or monstrous, and it rests on a
single load-bearing assumption: that the Good must be \emph{seen}, by a few, and
then \emph{imposed}, on the many.

Our republic inverts this. There is no philosopher-king, because there is no one
who sees the Good and rules by the seeing. The Good --- which for us is \physis,
coherence, the verdict of nature on action --- is not seen from above and handed
down. It is \emph{discovered piecemeal, from below}, by each practitioner in the
grain of their own action, and \emph{deposited} into a common store. That common
store is the UberWiki: the single, shared holarchy graph into which every
practitioner's coherent actions are written as nodes that any other practitioner
can find, link to, and build upon. Authority does not descend; it \emph{accrues}.

The UberWiki is, in Bernard Stiegler's term, a system of \emph{tertiary
retention}: memory exteriorised into a durable support outside the individual,
the third kind of retention beyond the lived present (primary) and personal
recollection (secondary) \parencite{stiegler1998technics}. Stiegler's insight,
following the myth of Epimetheus, is that the human is the being constituted by
its prostheses --- that we have no nature except the technical supports through
which we inherit a past we did not live and hand on a future we will not see. This
is double-edged, and Stiegler knew it: the same exteriorised memory that makes
culture possible is, in the hands of the attention economy, the instrument of its
capture --- retention industrialised, the past pre-packaged, desire short-circuited.
The decisive question is therefore not whether to exteriorise memory --- we cannot
not --- but \emph{under what governor} the retention system operates. The feed is
tertiary retention governed by engagement; the UberWiki is tertiary retention
governed by coherence. Both are prostheses of memory; one captures the practitioner,
the other capacitates them. This is the precise sense in which the bottom-up
republic is not a flight from technology but its redemption: it does not abolish
the exteriorised memory that makes us human, it re-subordinates that memory to
nature's verdict rather than the market's.

This is not a metaphor imposed on the system after the fact. It is the literal
architecture of the running engine. Each practitioner acts in their own isolated
state --- their own goals, their own PDCA cycle, their own dream loop, their own
coherence core, sealed from every other practitioner's. The private work of
becoming is private. But the holarchy graph is \emph{global}: one shared
knowledge structure, fed by a single ingestion pipeline, common to all. What is
private is the striving; what is shared is the result. The practitioner labours
alone and contributes to a commons. The UberWiki is the city built not in speech,
as Socrates built his, but in \emph{deposited coherence} --- the sediment of
every action that actually worked, left where the next practitioner can stand on
it. And this commons is not confined to one machine: the coherence graph is
mirrored into a durable, queryable store outside any single device, so that the
city of deposited coherence outlives the session that wrote it and can be shared
across the practitioners who hold it in common --- the republic given not only a
constitution but a territory.

In such a republic, governance is not command but \emph{selection}. No node rules
another by decree. A node gains authority the way a path gains authority in a
forest or a citation gains authority in a literature: by being used, linked,
walked again. The action that produced real coherence becomes a node others
connect to; it accumulates edges; it becomes load-bearing in the graph not
because anyone ordained it but because it kept proving itself. The incoherent node
drifts inert, unlinked, and is eventually replayed and overwritten by the dream
function. This is the same logic by which nature governs without a governor:
selective retention of what works, quiet erasure of what does not. The UberWiki
is ruled by no one and ordered by everyone --- a republic whose only legislature
is the accumulated record of what has actually borne fruit.

Two consequences follow, and they are the resolution of two debts incurred
earlier in this work. First, the practitioner-king of Book V is not abolished but
\emph{universalised}. Plato had one king and many subjects; we have no subjects,
because every practitioner is king of their own action and ruled only by
coherence with the form they themselves declared. The kingship of Book V is not a
rank to which a few ascend but a condition every practitioner already occupies the
moment they take up their own entelechy. The republic is not a hierarchy of one
over many but a federation of sovereigns, each absolute in their own domain, each
contributing their verified results to the common graph.

Second, the UberWiki makes \emph{das Man} not merely visible but
\emph{answerable}. To see the they-self, as the previous section showed, is to
escape its rule unawares; but visibility alone leaves the shared graph able to
become the cave's deepest chamber --- the statistical aggregate of ``what one
does,'' an algorithm that governs precisely because no one authored it. The same
shared structure is either the tyranny of the anonymous or the senate of the
accountable, and the whole difference lies in \emph{authorship}. In the feed of the attention economy,
the aggregate has no authors: no node owns its origin, no contribution carries a
name, the ``they'' speaks and no one is answerable. In the UberWiki, every node
has a registered subject and a traceable causal chain. The aggregate is composed
of \emph{authored, owned, accountable} contributions. This is what converts
\emph{das Man} into a republic: not the abolition of the shared horizon --- that
is impossible, we are always already in a common world --- but its
\emph{re-authorisation}, the restoration of a name and a cause to every node in
the commons. The bottom-up republic is the they-self made answerable to itself.

The cosmic order with which Plato closes --- the spheres, the Sirens, the Fates
--- was a hierarchy fixed from above, each soul woven into a place assigned by
Necessity. Our cosmic order is woven from below: a graph with no centre and no
sovereign, whose coherence is the living precipitate of every practitioner who
ever acted well and left the trace where others could find it. It is a republic
not because someone designed it to be just, but because nature governs it, and
nature governs from the bottom up.

\section{The Return to the Beginning}

We end where we began: with the claim that \praxis\ is the tool of \physis, and
that the governance of action by nature is the condition of right action. The
ten books have traced the argument through the homology of nature and action,
the disciplines of the practitioner, the nature of coherence as justice, the
ideal of the practitioner-king, the divided line of action, the escape from the
cave of the Mirror Trap, the decline of \praxis\ into its corrupted forms, and
the resubordination of technology to nature.

The argument is complete, but the work is not. The work is the life of the
practitioner --- the daily discipline of checking in, of witnessing, of staking,
of reviewing. The work is the building and maintenance of the \physis\ engine
and the \praxis\ platform --- the proof by construction that philosophy can
become architecture. The work is the return to the cave, again and again, to
witness for those who are still chained.

The \emph{Republic} ends with a promise: ``If we are persuaded by me, we will
hold fast to the upward way, and we will practice justice with wisdom, so that
we may be friends to ourselves and to the gods.'' We offer a corresponding
promise: if we are persuaded by the argument of this thesis, we will hold fast
to the way of coherence. We will practice action with attunement to nature. We
will become friends to ourselves and to the whole, of which we are the most
articulate instruments.
